\documentclass[11pt]{amsart}
\usepackage{amsmath,amssymb,amsthm}
\usepackage[margin=1.1in]{geometry}
\usepackage{booktabs,longtable}
\usepackage[colorlinks=true,allcolors=blue]{hyperref}
\usepackage[capitalise,nameinlink]{cleveref}

\newtheorem{theorem}{Theorem}[section]
\newtheorem{proposition}[theorem]{Proposition}
\newtheorem{lemma}[theorem]{Lemma}
\newtheorem{corollary}[theorem]{Corollary}
\theoremstyle{definition}
\newtheorem{definition}[theorem]{Definition}
\theoremstyle{remark}
\newtheorem{remark}[theorem]{Remark}

\newcommand{\E}{\mathbb{E}}
\newcommand{\PP}{\mathbb{P}}
\newcommand{\R}{\mathbb{R}}
\newcommand{\one}{\mathbf{1}}
\newcommand{\Simp}{\Delta_3}

\title[The three-state Potts model on regular trees]{The Kesten--Stigum bound is tight for the three-state Potts model on every regular tree}
\author{Yitzchak Shmalo}
\date{\today}

\begin{document}

\begin{abstract}
A uniformly random colour from $\{1,2,3\}$ is broadcast from the root of the infinite $d$-ary tree; each
edge keeps its parent's colour with probability $\lambda+(1-\lambda)/3$ and otherwise moves to one of the
two other colours uniformly. We prove that for every $d\ge2$ and every $0\le\lambda\le d^{-1/2}$ the
colours at depth $n$ carry asymptotically no information about the root. Together with the theorem of
Kesten and Stigum this shows that the reconstruction threshold of the ferromagnetic three-state Potts
channel on the $d$-ary tree is $\lambda=d^{-1/2}$ for every degree. Sly, and Mossel, Sly and Sohn, proved
this for all degrees above a constant that was not made explicit. Our proof makes that constant explicit,
$d\ge91$, by a quantitative Gaussian approximation, and covers $2\le d\le90$ by computer-assisted
arguments: an exact finite-state upgrade of the belief-propagation recursion bounds the second moment of
the root posterior at a finite depth, and exact polynomial certificates exclude every nontrivial
belief-propagation fixed point below that bound.
\end{abstract}

\maketitle

\section{Introduction}

Let $T_d$ be the infinite rooted $d$-ary tree, $q\ge2$ and $\lambda\in[0,1]$, and let
$M_\lambda=\lambda I+(1-\lambda)J$ be the $q\times q$ Potts channel, where every entry of $J$ is $1/q$.
The broadcast process gives the root $\rho$ a uniform colour $\sigma_\rho\in[q]$ and passes colours
down each edge independently through $M_\lambda$. Write $\sigma_{L_n}$ for the colours at depth $n$. The
process is \emph{reconstructible} if
\[
\lim_{n\to\infty}\E\,\big\|\PP(\sigma_\rho=\cdot\mid\sigma_{L_n})-q^{-1}\one\big\|_{\mathrm{TV}}>0
\]
and \emph{nonreconstructible} otherwise. Kesten and Stigum~\cite{KS66} showed that reconstruction holds
whenever $d\lambda^2>1$. For $q=2$ the bound is sharp~\cite{BRZ95,EKPS00}. For large $q$ it is not:
Mossel~\cite{Mos01} showed that reconstruction can hold below it, and Sly~\cite{Sly11} proved that the
bound fails to be tight for $q\ge5$ and that it is tight for $q=3$ once $d\ge d_0$. Mossel, Sly and
Sohn~\cite{MSS22} extended tightness to $q=4$ and to Galton--Watson trees of large mean. In both papers
$d_0$ comes from a central limit approximation and is not computed, and Mossel, Sly and Sohn ask whether
the result holds for $q=3$ and small $d$. Bhatnagar and Maneva~\cite{BM11} gave a computational method
that bounds reconstruction thresholds on small trees; a finite computation of this kind cannot by itself
reach the Kesten--Stigum bound, since at $d\lambda^2=1$ the information decays only polynomially in the
depth.

\begin{theorem}\label{thm:main}
Let $q=3$. For every integer $d\ge2$ and every $0\le\lambda\le d^{-1/2}$, the broadcast process on $T_d$
is nonreconstructible. Hence the reconstruction threshold of the ferromagnetic three-state Potts channel
on the $d$-ary tree equals the Kesten--Stigum bound $\lambda=d^{-1/2}$ for every $d\ge2$.
\end{theorem}

The proof follows one plan at every degree. Let $x_n=\E\sum_i(p_i-1/3)^2$, where
$p=\PP(\sigma_\rho=\cdot\mid\sigma_{L_n})$. The sequence $x_n$ is nonincreasing, its limit $x_\infty$
is the second moment of a colour-symmetric fixed point of the belief-propagation map
(\cref{sec:fixed}), and nonreconstruction is the statement $x_\infty=0$. Two estimates meet:
\begin{enumerate}
\item an \emph{exclusion radius} $r_d>0$: no nontrivial symmetric fixed point at the critical
  $\lambda=d^{-1/2}$ has second moment in $(0,r_d]$ (\cref{sec:local});
\item an \emph{entry bound}: $x_N<r_d$ for some explicit depth $N$ (\cref{sec:envelope,sec:large}).
\end{enumerate}
At the critical point the drift of $x$ near zero is $-\tfrac34\lambda^4d(d-1)x^2$ to leading order
(\cite[(4.9)]{Sly11}), so the fixed-point equation leaves no room for a small nonzero solution; the work in
(1) is to make the neighbourhood explicit. We do this with exact Bernstein-basis certificates for
polynomial inequalities in the low moments of the fixed point. For (2) at $2\le d\le90$ we run the
belief-propagation recursion on an information upgrade of the posterior law that is supported on a
finite grid, in exact integer arithmetic; the upgrade can only increase $x_n$, so the computed value is an
upper bound. For $d\ge91$ we replace the computation by a Lindeberg estimate with every constant written
down, comparing the recursion with the Gaussian map of Sly~\cite{Sly11}, together with a certified lower
bound on the gap between that map and the identity. Monotonicity in $\lambda$ reduces subcritical
$\lambda$ to the critical one (\cref{sec:proof}).

Table~\ref{tab:degrees} in \cref{sec:proof} lists, for each degree or band of degrees, the radius and the
certified bound. The computations are in the accompanying repository and are described in
\cref{sec:computations}.

\section{The recursion and its fixed points}\label{sec:fixed}

Throughout, $q=3$, $u=(\tfrac13,\tfrac13,\tfrac13)$, and $\Simp$ is the probability simplex in $\R^3$ with
the Euclidean norm $|\cdot|$. The symmetric group $S_3$ acts on $\Simp$ by permuting coordinates. For a
Borel probability measure $\mu$ on $\Simp$ write
\[
x(\mu)=\int|p-u|^2\,\mu(dp).
\]
We say that $\mu$ is \emph{symmetric} if it is $S_3$-invariant, and that $\mu$ is a \emph{posterior law} if
there are a random colour $\sigma$, uniform on $[3]$, and a $\sigma$-algebra $\mathcal G$ such that $\mu$ is
the law of $\PP(\sigma=\cdot\mid\mathcal G)$. A posterior law has mean $u$, and $x(\mu)=0$ only for the
point mass $\delta_u$.

\subsection{The recursion}
Fix $d\ge2$ and $0\le\lambda<1$. For a probability measure $\mu$ on $\Simp$ with mean $u$ let
$p^{(1)},\dots,p^{(d)}$ be independent with law $\mu$, and put
\[
z_i=\prod_{j=1}^d\big(1-\lambda+3\lambda p^{(j)}_i\big)\quad(i=1,2,3),\qquad S=z_1+z_2+z_3 .
\]
Define $T_\lambda\mu$ by
\begin{equation}\label{eq:T}
\int h\,d(T_\lambda\mu)=\E\Big[h\big(z/S\big)\,\frac S3\Big]
\end{equation}
for bounded measurable $h$. Each factor $1-\lambda+3\lambda p_i$ has mean $1$ under $\mu$, so $\E z_i=1$
and $\E S/3=1$; thus $T_\lambda\mu$ is a probability measure, and it again has mean $u$ because
$\int p_i\,d(T_\lambda\mu)=\E[z_i]/3=1/3$.

Let $P_n=\PP(\sigma_\rho=\cdot\mid\sigma_{L_n})$ and let $\mu_n$ be its law. Thus $\mu_0$ is the uniform law
on the three vertices of $\Simp$.

\begin{lemma}\label{lem:recursion}
$\mu_{n+1}=T_\lambda\mu_n$ for every $n\ge0$. Each $\mu_n$ is a symmetric posterior law.
\end{lemma}

\begin{proof}
Let $v_1,\dots,v_d$ be the children of $\rho$, and let $Y_j$ be the colours at depth $n+1$ below $v_j$, so
that $\sigma_{L_{n+1}}=(Y_1,\dots,Y_d)$. The subtree below $v_j$ is a copy of $T_d$, and $\sigma_{v_j}$ is
uniform; hence $Q^{(j)}=\PP(\sigma_{v_j}=\cdot\mid Y_j)$ has law $\mu_n$. The $Y_j$ are not independent,
since they share the root; we compute their joint law. By Bayes' rule with a uniform prior,
$\PP(Y_j\mid\sigma_{v_j}=k)=3Q^{(j)}_k\,\PP(Y_j)$. Summing over the colour of $v_j$,
\[
\PP(Y_j\mid\sigma_\rho=i)=\sum_k M_\lambda(i,k)\,3Q^{(j)}_k\,\PP(Y_j)
=\big(1-\lambda+3\lambda Q^{(j)}_i\big)\PP(Y_j).
\]
Since the subtrees are conditionally independent given $\sigma_\rho$,
$\PP(Y_1,\dots,Y_d\mid\sigma_\rho=i)=z_i\prod_j\PP(Y_j)$ with $z$ built from the $Q^{(j)}$. The posterior of
the root is therefore $z/S$, and, averaging over the uniform root, the joint law of $(Y_1,\dots,Y_d)$ has
density $\tfrac13\sum_i z_i=S/3$ with respect to the product of its marginals. Since $z$ is a function of
$(Q^{(1)},\dots,Q^{(d)})$, the vector $(Q^{(1)},\dots,Q^{(d)})$ has density $S/3$ with respect to
$\mu_n^{\otimes d}$, which is~\eqref{eq:T}. Symmetry passes from $\mu_n$ to $\mu_{n+1}$ because
$M_\lambda$ commutes with permutations of the colours, and $\mu_n$ is a posterior law by definition.
\end{proof}

\subsection{The tail limit}
Let $\mathcal F_n$ be the $\sigma$-algebra generated by all colours at depth at least $n$. These
$\sigma$-algebras decrease. Given $\sigma_{L_n}$, the root is independent of the colours below depth $n$, so
$P_n=\PP(\sigma_\rho=\cdot\mid\mathcal F_n)$. Put $x_n=x(\mu_n)=\E|P_n-u|^2$ and
$\mathcal F_\infty=\bigcap_n\mathcal F_n$.

\begin{lemma}\label{lem:tail}
The sequence $x_n$ is nonincreasing. The posteriors $P_n$ converge almost surely and in $L^2$ to
$P_\infty=\PP(\sigma_\rho=\cdot\mid\mathcal F_\infty)$, whose law $\mu_\infty$ is a symmetric posterior law
with $T_\lambda\mu_\infty=\mu_\infty$ and $x(\mu_\infty)=\lim_n x_n\le x_N$ for every $N$.
\end{lemma}

\begin{proof}
$(P_n)$ is a reverse martingale for the decreasing filtration $(\mathcal F_n)$, so
$P_{n+1}=\E[P_n\mid\mathcal F_{n+1}]$, and Jensen's inequality gives $x_{n+1}\le x_n$. L\'evy's downward
theorem gives convergence almost surely, and in $L^2$ by boundedness, to $P_\infty$. Hence $\mu_n\to\mu_\infty$
weakly and $x_n\to x(\mu_\infty)$. Symmetry passes to the limit, and $\mu_\infty$ is a posterior law by
construction.

For the fixed-point identity, each factor in $z$ is at least $1-\lambda>0$, so $S\ge3(1-\lambda)^d$, and
the integrand in~\eqref{eq:T} is a bounded continuous function of $(p^{(1)},\dots,p^{(d)})\in\Simp^d$ when
$h$ is bounded and continuous. Weak convergence $\mu_n\to\mu_\infty$ implies weak convergence of the product
laws $\mu_n^{\otimes d}\to\mu_\infty^{\otimes d}$, hence $T_\lambda\mu_n\to T_\lambda\mu_\infty$ weakly. By
\cref{lem:recursion}, $T_\lambda\mu_n=\mu_{n+1}\to\mu_\infty$, so $T_\lambda\mu_\infty=\mu_\infty$.
\end{proof}

\begin{proposition}\label{prop:reduction}
Let $0<\lambda<1$ and $r>0$. Suppose that every symmetric posterior law $\mu$ with $T_\lambda\mu=\mu$ and
$x(\mu)\le r$ equals $\delta_u$, and that $x_N\le r$ for some $N$. Then the broadcast process with
parameters $(d,\lambda)$ is nonreconstructible.
\end{proposition}

\begin{proof}
By \cref{lem:tail}, $\mu_\infty$ satisfies the hypotheses, so $\mu_\infty=\delta_u$ and $x_n\to0$. By
Cauchy--Schwarz in $\R^3$ and Jensen,
$\E\|P_n-u\|_{\mathrm{TV}}=\tfrac12\E\sum_i|P_n(i)-\tfrac13|\le\tfrac{\sqrt3}2\sqrt{x_n}\to0$.
\end{proof}

\subsection{Comparison of experiments}
Write $x_n(d,\lambda)$ for $x_n$ when the parameters need to be named.

\begin{lemma}\label{lem:compare}
If $2\le d\le D$ and $0\le\lambda\le\eta<1$, then $x_n(d,\lambda)\le x_n(D,\eta)$ for every $n$.
\end{lemma}

\begin{proof}
Both changes are garblings of the observation that do not depend on the root. If $Y=g(X,U)$ with $U$
independent of $(\sigma_\rho,X)$, then $\PP(\sigma_\rho=\cdot\mid Y)=\E[\PP(\sigma_\rho=\cdot\mid X)\mid Y]$,
and Jensen's inequality for the convex function $p\mapsto|p-u|^2$ shows that the second moment can only
decrease. For the degree, $T_d$ is a subtree of $T_D$ containing the root, and the depth-$n$ colours of
$T_d$ are a subset of those of $T_D$. For the strength, if $\eta>0$ then
\[
M_\lambda=\tfrac\lambda\eta\,M_\eta+\big(1-\tfrac\lambda\eta\big)J,
\]
and an edge of strength $\lambda$ is an edge of strength $\eta$ with probability $\lambda/\eta$ and
otherwise gives the child a uniform colour independent of its parent. Mark each edge independently with
probability $1-\lambda/\eta$. On the vertices joined to the root by unmarked paths, keep the colours of the
strength-$\eta$ process; below every marked edge with no marked edge above it, put a fresh copy of the
strength-$\lambda$ process started from a uniform colour. The result is the strength-$\lambda$ process, and
its depth-$n$ colours are a function of the depth-$n$ colours of the strength-$\eta$ process together with
the marks and the fresh copies, which are independent of the root and of that process.
\end{proof}

By \cref{lem:compare} with $D=d$, nonreconstruction at $\lambda=d^{-1/2}$ implies it for every smaller
$\lambda$, and an upper bound on $x_N(D,\eta)$ with $D\ge d$ and $\eta\ge d^{-1/2}$ is an upper bound on
$x_N(d,d^{-1/2})$. The second use lets one computation serve a band of degrees (\cref{sec:envelope}).

\section{Exclusion near the uniform law}\label{sec:local}

In this section $\lambda=d^{-1/2}$ unless stated otherwise. We prove that a nontrivial symmetric fixed point
cannot have small second moment, with the radii $r_2=\tfrac{23}{2000}$, $r_3=\tfrac{3}{100}$ and, for $d\ge4$,
\begin{equation}\label{eq:radii}
r_d=\max\Big\{\tfrac5{36}\big(1-d^{-1/2}\big)^2,\ \varrho_d\Big\},\qquad
\varrho_d=\begin{cases}0,&d=4,\\ \tfrac{217}{4500},&5\le d\le19,\\ \tfrac{257}{3000},&20\le d\le79,\\
\tfrac{179}{1500},&d\ge80.\end{cases}
\end{equation}

\subsection{Fusion and moments}
For $p\in\Simp$ write $a(p)=3p-\one$, a vector with coordinate sum zero and $|a|^2\le6$. For a symmetric
probability measure $\mu$ on $\Simp$ put
\[
M(\mu)=\int|a|^2d\mu,\qquad U(\mu)=\int|a|^4d\mu,\qquad T(\mu)=\int\textstyle\sum_ia_i^3\,d\mu,
\]
so that $x(\mu)=M(\mu)/9$ and $0\le M\le6$. For symmetric $\mu,\nu$ define the \emph{fusion} $\mu\circledast\nu$ by
\[
\int h\,d(\mu\circledast\nu)=\E\Big[h\Big(\frac{p\odot r}{\langle p,r\rangle}\Big)\,3\langle p,r\rangle\Big],
\qquad p\sim\mu,\ r\sim\nu\ \text{independent},
\]
where $\odot$ is the coordinatewise product. If $\mu$ and $\nu$ are the laws of the posteriors of a uniform colour
$\sigma$ given two observations that are conditionally independent given $\sigma$, then, exactly as in the proof
of \cref{lem:recursion}, $\mu\circledast\nu$ is the law of the posterior given both. Hence $\circledast$ is
commutative and associative on symmetric posterior laws, and $M(\mu\circledast\nu)\ge M(\mu)$ by the argument of
\cref{lem:compare}. Let $E_\lambda\mu$ be the law of $u+\lambda(p-u)$ for $p\sim\mu$; it is the law of the
posterior of $\rho$ given the subtree of one child whose own posterior has law $\mu$. Since
$1-\lambda+3\lambda p_i=3(u+\lambda(p-u))_i$,
\[
T_\lambda\mu=(E_\lambda\mu)^{\circledast d},\qquad
M(E_\lambda\mu)=\lambda^2M(\mu),\quad U(E_\lambda\mu)=\lambda^4U(\mu),\quad T(E_\lambda\mu)=\lambda^3T(\mu).
\]
Every $p$ in the support of $E_\lambda\mu$ has $p_i\ge(1-\lambda)/3$.

For $p,r\in\Simp$ put $\alpha=a(p)$, $\beta=a(r)$, $c=\langle\alpha,\beta\rangle$, $S=3+c=9\langle p,r\rangle$ and
\[
v=\alpha+\beta+\alpha\odot\beta-\tfrac c3\one,\qquad N=|v|^2 .
\]
A direct computation gives $a\big(p\odot r/\langle p,r\rangle\big)=3v/S$, so the fused vector has $a$-norm
$9N/S^2$ and weight $S/3$.

\begin{lemma}\label{lem:variance}
Let $\mu,\nu$ be symmetric probability measures on $\Simp$ with moments $(m,V,T_\alpha)$ and $(n,W,T_\beta)$
respectively, and suppose $S\ge s>0$ on the support of $\mu\otimes\nu$. Then
\[
M(\mu\circledast\nu)=m+n-\tfrac{mn}6-\tfrac13(nT_\alpha+mT_\beta)+R,\qquad 0\le R\le\frac{\E[Nc^2]}{3s},
\]
\[
\E[Nc^2]=\tfrac12(nV+mW)+3T_\alpha T_\beta+\tfrac12(VT_\beta+T_\alpha W)+\tfrac1{12}VW .
\]
\end{lemma}

\begin{proof}
By the remark before the lemma, $M(\mu\circledast\nu)=\E[(9N/S^2)(S/3)]=3\E[N/S]$. Since
$(\tfrac13-\tfrac c9)S=1-\tfrac{c^2}9$, we have $1/S=\tfrac13-\tfrac c9+\tfrac{c^2}{9S}$, hence
$M(\mu\circledast\nu)=\E N-\tfrac13\E[Nc]+R$ with $R=\E[Nc^2/(3S)]$, which lies between $0$ and
$\E[Nc^2]/(3s)$. A symmetric law of a zero-sum vector $\alpha$ with $\E|\alpha|^2=m$,
$\E|\alpha|^4=V$, $\E\sum\alpha_i^3=T_\alpha$ has $\E\alpha=0$ and, for distinct indices $i,j,k$,
\[
\E\alpha_i^2=\tfrac m3,\ \E\alpha_i\alpha_j=-\tfrac m6;\quad
\E\alpha_i^3=\tfrac{T_\alpha}3,\ \E\alpha_i^2\alpha_j=-\tfrac{T_\alpha}6,\ \E\alpha_1\alpha_2\alpha_3=\tfrac{T_\alpha}3;
\]
\[
\E\alpha_i^4=\tfrac V6,\ \E\alpha_i^3\alpha_j=-\tfrac V{12},\ \E\alpha_i^2\alpha_j^2=\tfrac V{12},\
\E\alpha_i^2\alpha_j\alpha_k=0 .
\]
(The second and fourth tables follow from symmetry and $\sum_i\alpha_i=0$; for the fourth, also
$\sum_i\alpha_i^4=\tfrac12|\alpha|^4$ for zero-sum triples.) Expanding $N$, $Nc$ and $Nc^2$ as polynomials in
$\alpha,\beta$ and contracting with these tables gives $\E N=m+n+\tfrac{mn}6$,
$\E[Nc]=mn+nT_\alpha+mT_\beta$ and the stated formula for $\E[Nc^2]$.
\end{proof}

\subsection{Three moment inequalities}
Fix $\eta=2887/5000$ and $\kappa=2(99/140)^4<\tfrac12+10^{-4}$.

\begin{lemma}\label{lem:moments}
Let $\mu,\nu$ be symmetric posterior laws with moments $(M_\mu,U_\mu,T_\mu)$ and $(M_\nu,U_\nu,T_\nu)$.
\begin{enumerate}
\item[\textup{(CF)}] $T(\mu\circledast\nu)\ge(1-\tfrac{M_\nu}6)T_\mu+(1-\tfrac{M_\mu}6)T_\nu$.
\item[\textup{(F)}] If $\nu$ is supported on $\{r:r_i\ge(1-\eta)/3\}$, then
$U(\mu\circledast\nu)\le U_\mu+U_\nu+5M_\mu M_\nu$.
\item[\textup{(C)}] Under the same hypothesis, $|T(\mu\circledast\nu)-T_\mu-T_\nu|\le\tfrac53M_\mu M_\nu$.
\item[\textup{(F1)}] If $0<\lambda\le1/\sqrt2$, then
$\lambda^4U\big((E_\lambda\mu)^{\circledast2}\big)\le\kappa\lambda^4U_\mu+\tfrac32\lambda^4M_\mu^2$.
\end{enumerate}
\end{lemma}

\begin{proof}
Each inequality is bilinear in the pair $(\mu,\nu)$ once the terms that involve one law only are multiplied by
the total mass of the other, and the fusion is bilinear by its definition. A symmetric law is a mixture of
uniform laws on $S_3$-orbits, so it suffices to prove each inequality when $\mu$ and $\nu$ are uniform on the
orbits of two points. Write the sorted points as
\[
p(t,y)=\big(\tfrac{1-t}3,\ \tfrac{1-t}3+\tfrac{ty}2,\ \tfrac{1-t}3+t(1-\tfrac y2)\big),\qquad 0\le t,y\le1,
\]
and similarly $r=p(s,z)$. The fused law is then the average over the six relative permutations $\pi$ of $r$, the
term for $\pi$ carrying weight $S_\pi/18$ and output $a$-vector $3v_\pi/S_\pi$. After clearing the positive
denominators $S_\pi$ (in (CF) the terms with $S_\pi=0$ carry no weight), each inequality becomes the
nonnegativity of an explicit polynomial with integer coefficients on a product of intervals in $(t,y,s,z)$.
We prove this nonnegativity in \cref{sec:computations} by dividing out the monomial factors that vanish on the
boundary and expanding the quotient in a Bernstein basis, all of whose coefficients are nonnegative. For (F1)
the reduction is first made at $\lambda=99/140$; the case of smaller $\lambda$ follows by \cref{lem:compare}
applied to the two-child tree, since the right side is then computed with the same scale.
\end{proof}

\begin{corollary}\label{cor:sign}
If $d\lambda^3<1$, then every symmetric posterior law $\mu$ with $T_\lambda\mu=\mu$ has $T(\mu)\ge0$.
\end{corollary}

\begin{proof}
Let $T=T(\mu)$ and let $T_j$ be the cubic moment of $(E_\lambda\mu)^{\circledast j}$, so $T_1=\lambda^3T$ and
$T_d=T$. Suppose $T<0$. Since the coefficients in (CF) lie in $[0,1]$, (CF) gives
$T_{j+1}\ge\min(T_j,0)+\lambda^3T$, hence $T_j\ge j\lambda^3T$ for all $j$. At $j=d$ this reads
$(1-d\lambda^3)T\ge0$, a contradiction.
\end{proof}

\subsection{Degree two}

\begin{proposition}\label{prop:local2}
Let $d=2$ and $\lambda=2^{-1/2}$. A symmetric posterior law $\mu$ with $T_\lambda\mu=\mu$ and $x(\mu)\le r_2$
equals $\delta_u$.
\end{proposition}

\begin{proof}
Let $m,V,t$ be the moments of the child law $E_\lambda\mu$, so $m=M(\mu)/2=\tfrac92x(\mu)$, $V=U(\mu)/4$ and
$\mu=(E_\lambda\mu)^{\circledast2}$. By (F1) at the fixed point, $V\le\kappa V+\tfrac32m^2$, so
$V\le Cm^2$ with $C=\tfrac{3001}{1000}$. For a zero-sum vector, $|\sum_i\alpha_i^3|\le|\alpha|^3/\sqrt6$, so by
Cauchy--Schwarz $|t|\le\sqrt{mV/6}\le\theta m^{3/2}$ with $\theta=\tfrac{177}{250}\ge\sqrt{C/6}$. For two child
vectors $c\ge-3\lambda^2$, so $S\ge3/2$. \Cref{lem:variance} with both laws equal to $E_\lambda\mu$ and the
identity $M(\mu)=2m$ give
\[
\tfrac{m^2}6\le\tfrac23m|t|+\tfrac29\big(mV+3t^2+V|t|+\tfrac1{12}V^2\big).
\]
If $m>0$, divide by $m^2$ and insert the bounds above:
\[
\tfrac16\le\tfrac23\theta\sqrt m+\tfrac29\big(\tfrac32Cm+C\theta m^{3/2}+\tfrac1{12}C^2m^2\big).
\]
For $x(\mu)\le\tfrac{23}{2000}$ we have $m\le\tfrac{207}{4000}$ and $\sqrt m\le\tfrac{91}{400}$, and the right side
is at most $0.16516<\tfrac16$ (exact rational arithmetic). Hence $m=0$.
\end{proof}

\subsection{Degree at least three}
Let $d\ge3$, $\lambda=d^{-1/2}$, and let $\mu$ be a nontrivial symmetric posterior fixed point with moments
$M,U,T$. Let $\mu_j=(E_\lambda\mu)^{\circledast j}$, with moments $M_j,U_j,T_j$, so $\mu_d=\mu$ and
$M_j\le M$. The child law has moments $M/d$, $U/d^2$, $\lambda T/d$, and its support satisfies $r_i\ge(1-\lambda)/3$
with $\lambda<\eta$. By \cref{cor:sign} (note $d\lambda^3=\lambda<1$) and (CF), $T\ge0$ and $T_j\ge0$ for all $j$.
For $2\le k\le d$ set
\[
g_k=\frac{k(k-1)}{12d^2},\qquad Z=\frac{k(k-1)}2+d(d-k),\qquad C_U=\frac{5Z}{d(d-1)},\qquad
C_T=\frac{C_U(1+\lambda)}3,
\]
and let $A_k,B_k\ge0$ be the closed forms
\[
A_k=A_2+\frac{S_2}{72d^3}+\frac{K(C_US_1+\tfrac54S_2)}{d^3},\qquad
B_k=B_2+K\Big[\big(3t_e+\tfrac{v_e}2\big)\Sigma_t+\big(\tfrac{t_e}2+\tfrac{v_e}{12}\big)\Sigma_u\Big],
\]
where $S_1=\tfrac{k(k-1)}2-1$, $S_2=\tfrac{k(k-1)(k-2)}3$, $K=\frac1{9(1-\lambda)}$, $K'=\frac d{9(d-1)}$,
$v_e=C_U/d^2$, $t_e=\lambda C_T/d$, $\Sigma_u=(C_US_1+\tfrac52S_2)/d^2$, $\Sigma_t=\lambda C_TS_1/d+\tfrac56S_2/d^2$,
$A_2=K'C_U/d^3$ and $B_2=K'\big[3C_T^2/d^3+\lambda C_UC_T/d^3+C_U^2/(12d^4)\big]$.

\begin{lemma}\label{lem:prefix}
Let $0<M_*\le6$ and $0<\delta<1$, and suppose that
\begin{equation}\label{eq:criterion}
(1-\delta)g_k-M_*A_k-M_*^2B_k\ge0\qquad(2\le k\le d).
\end{equation}
Then no nontrivial symmetric posterior fixed point has $M\le M_*$.
\end{lemma}

\begin{proof}
Suppose $0<M\le M_*$. We show by induction on $k=2,\dots,d$ that
\begin{equation}\label{eq:prefix}
M_k\le\tfrac kdM-g_kM^2+A_kM^3+B_kM^4,
\end{equation}
and hence, by \eqref{eq:criterion}, $M_k\le\tfrac kdM-\delta g_kM^2<\tfrac kdM$. At $k=d$ this contradicts
$M_d=M$.

Assume $M_j\le jM/d$ for $1\le j<k$ (an equality at $j=1$). Together with $M_j\le M$ for $j\ge k$ this gives
$\sum_{j=1}^{d-1}M_j\le M\big[\tfrac{k(k-1)}{2d}+d-k\big]=MZ/d$. Applying (F) and (C) to
$\mu_{j+1}=\mu_j\circledast E_\lambda\mu$ and telescoping from $j=1$ to $j=d-1$,
\[
\big(1-\tfrac1d\big)U\le\tfrac{5M}d\sum_{j<d}M_j,\qquad(1-\lambda)T\le\tfrac{5M}{3d}\sum_{j<d}M_j,
\]
so, using $(1-\lambda)(1+\lambda)=(d-1)/d$, $U\le C_UM^2$ and $0\le T\le C_TM^2$. Telescoping the same
inequalities only up to $j<k$ gives $U_j\le u_jM^2$ and $T_j\le t_jM^2$ with
$u_j=[jC_U+\tfrac52j(j-1)]/d^2$ and $t_j=j\lambda C_T/d+\tfrac56j(j-1)/d^2$.

For $k=2$, apply \cref{lem:variance} to two child laws: their vectors satisfy $c\ge-3\lambda^2$, so
$S\ge3(d-1)/d$, and the cubic term is nonpositive since $T\ge0$. This gives \eqref{eq:prefix} at $k=2$. For the
step from $j$ to $j+1<k+1$, apply \cref{lem:variance} with left law $\mu_j$ and right law $E_\lambda\mu$: here
$S\ge3(1-\lambda)$ because every right coordinate is at least $(1-\lambda)/3$, the cubic term is nonpositive, and
the bounds above give
\[
R\le K\Big[\tfrac{u_j}{2d}+\tfrac{jC_U}{2d^3}\Big]M^3+K\Big[3t_jt_e+\tfrac12(u_jt_e+t_jv_e)+\tfrac1{12}u_jv_e\Big]M^4 .
\]
The main term is $(1-\tfrac M{6d})M_j+\tfrac Md$, with a positive multiplier since $M\le6$. Inserting
\eqref{eq:prefix} for $M_j$ and dropping the nonpositive products of the multiplier's $-M/(6d)$ with the cubic and
quartic terms, the quadratic coefficient becomes $g_j+j/(6d^2)=g_{j+1}$, and the cubic and quartic coefficients
grow by exactly the amounts whose sums over $j$ are the closed forms $A_k,B_k$ above.
\end{proof}

\begin{proposition}\label{prop:local3}
Let $d\ge3$ and $\lambda=d^{-1/2}$. A symmetric posterior law $\mu$ with $T_\lambda\mu=\mu$ and $x(\mu)\le r_d$
equals $\delta_u$.
\end{proposition}

\begin{proof}
We verify \eqref{eq:criterion} for all integers $d\ge3$ and $2\le k\le d$ at once. Substitute $\ell=\lambda$,
$d=\ell^{-2}$, $k=2+\rho(d-2)$ with $\rho\in[0,1]$, and $\ell=Lz$ with $z\in[0,1]$. After multiplication by a
positive factor $D(\ell)$ the left side of \eqref{eq:criterion} becomes a polynomial in $(z,\rho)$, and every
integer pair $(d,k)$ with $d^{-1/2}\le L$ corresponds to a point of $[0,1]^2$. The following choices are used.
\begin{center}
\begin{tabular}{lllll}
\hline
degrees & $M_*$ & $x=M_*/9$ & $\delta$ & $L$\\
\hline
$d\ge3$ & $27/100$ & $3/100$ & $1/100$ & $\eta=2887/5000$\\
$d\ge4$ & $\tfrac54(1-\ell)^2$ & $\tfrac5{36}(1-\lambda)^2$ & $1/1000$ & $1/2$\\
$d\ge5$ & $217/500$ & $217/4500$ & $1/1000$ & $559017/1250000$\\
$d\ge20$ & $771/1000$ & $257/3000$ & $1/1000$ & $559017/2500000$\\
$d\ge80$ & $537/500$ & $179/1500$ & $1/1000$ & $559017/5000000$\\
\hline
\end{tabular}
\end{center}
In each case $L\ge d^{-1/2}$ at the smallest degree of the row, and the Bernstein coefficients of the polynomial
on $[0,1]^2$ are nonnegative (\cref{sec:computations}). Hence \eqref{eq:criterion} holds and \cref{lem:prefix}
applies with $x=M/9$. For $d\ge4$ the curve dominates the constant $3/100$, since
$(1-d^{-1/2})^2\ge\tfrac14>\tfrac{27}{125}$, so $r_d$ in \eqref{eq:radii} is the largest radius available.
\end{proof}

\section{Certified upper envelopes}\label{sec:envelope}

By \cref{lem:tail}, an entry bound $x_N(d,d^{-1/2})<r_d$ is needed at one depth $N$ only. We obtain it from a
finite computation whose output is more informative than the true posterior law at every depth, so that its second
moment is an upper bound. The method is that of Bhatnagar and Maneva~\cite{BM11}. We state the version used here
in full, because the certificates depend on its details: observations of probability zero, the discretization
and the rounding of masses are all part of the argument.

\subsection{Two upgrades}
Let $\mu,\nu$ be symmetric posterior laws. Say that $\nu$ is \emph{more informative} than $\mu$ if there are
observations $X$ and $Y$ of a uniform colour $\sigma$, with posterior laws $\nu$ and $\mu$, such that
$Y=g(X,U)$ with $U$ independent of $(\sigma,X)$. As in the proof of \cref{lem:compare}, this implies
$x(\nu)\ge x(\mu)$. The relation is preserved by $E_\lambda$ and by fusion with any third law, since the garbling
can be applied to one observation and the others left alone.

Fix an integer $K\ge1$ and let $G_K=\{v/K:\ v\in\mathbb Z_{\ge0}^3,\ v_1+v_2+v_3=K\}$. For $p\in\Simp$ write
$Kp=f+r$ with $f$ integral and $0\le r_i<1$, so that $r_1+r_2+r_3\in\{0,1,2\}$. The \emph{split} of $p$ is the
point $p$ itself if this sum is $0$; the points $(f+e_i)/K$ with weights $r_i$ if it is $1$; and the points
$(f+\one-e_i)/K$ with weights $1-r_i$ if it is $2$. In each case the weights are nonnegative, sum to one and
average to $p$, and the split commutes with permutations of the colours.

\begin{lemma}\label{lem:split}
Let $\mu$ be a symmetric posterior law and let $\mu'$ be the law obtained by replacing each $p$ by its split.
Then $\mu'$ is a symmetric posterior law on $G_K$, more informative than $\mu$.
\end{lemma}

\begin{proof}
Let $X$ be an observation with posterior $p=\sum_jw_jq_j$, the split of $p$. Given $X$ and $\sigma=i$ with
$p_i>0$, draw a label $j$ with probability $w_j(q_j)_i/p_i$; these probabilities sum to one. If $p_i=0$ then
$(q_j)_i=0$ whenever $w_j>0$, the pair $(X,\sigma=i)$ has probability zero, and the kernel may be chosen
arbitrarily there. By Bayes' rule the posterior given $(X,j)$ is $q_j$, and given $X$ the label equals $j$ with
probability $w_j$. Forgetting the label recovers $X$. The posterior law of $(X,j)$ is $\mu'$.
\end{proof}

A symmetric law on $G_K$ is determined by the masses of its $S_3$-orbits; for $K=96$ there are $817$ orbits,
of sizes $1$, $3$ and $6$. Fix $S=2^b$. The \emph{rounding} of such a law replaces each orbit mass $w$ by
$\lfloor Sw\rfloor/S$ and adds the total removed mass to the orbit $\{e_1,e_2,e_3\}$.

\begin{lemma}\label{lem:round}
The rounding of a symmetric posterior law on $G_K$ is a symmetric posterior law, more informative than it.
\end{lemma}

\begin{proof}
For a posterior law, $\PP(X=p\mid\sigma=i)=3p_i\,\PP(X=p)$. The points of an orbit average to $u$, so the
orbit containing the observation is independent of $\sigma$. On an orbit of mass $w$, keep the observation
with probability $\lfloor Sw\rfloor/(Sw)$, independently of $\sigma$ and of the point within the orbit, and
otherwise reveal $\sigma$ in addition. The revealed observations have posterior $e_\sigma$, and $\sigma$ is
uniform on them because the kept fraction does not depend on $\sigma$. The refined observation has the rounded
law, and discarding the revelation recovers the original observation.
\end{proof}

Write $\Pi$ for a split followed by a rounding. By \cref{lem:split,lem:round}, $\Pi\mu$ is more informative
than $\mu$.

\subsection{The computed laws}
Fix $D\ge2$, a rational $\eta=A/B\in(0,1)$ in lowest terms, $K$ divisible by $3$, and $S=2^b$. Let $\nu_0$ be the
uniform law on $\{e_1,e_2,e_3\}$. Given $\nu=\nu_n$, the law $\nu_{n+1}$ is computed in one of two ways.

\emph{Sequential fusion} ($D\le5$). Put $P_2=\Pi\big(E_\eta\nu\circledast E_\eta\nu\big)$ and
$P_{j+1}=\Pi\big(P_j\circledast E_\eta\nu\big)$ for $2\le j<D$, and set $\nu_{n+1}=P_D$. The shrunken law
$E_\eta\nu$ enters each fusion exactly and is not split; every factor is a new copy of the child law $\nu$,
never the partial product $P_j$.

\emph{Binary fusion} ($D\ge6$). Put $\beta=\Pi(E_\eta\nu)$, $\beta_0=\beta$ and
$\beta_{k+1}=\Pi(\beta_k\circledast\beta_k)$, where the two factors are independent copies. If
$D=2^{k_1}+\dots+2^{k_s}$ with distinct $k_t$, the laws $\beta_{k_1},\dots,\beta_{k_s}$ are fused one at a time,
with $\Pi$ applied after each fusion, and the result is $\nu_{n+1}$.

\begin{proposition}\label{prop:dominate}
For every $n\ge0$, $\nu_n$ is more informative than $\mu_n(D,\eta)$, and hence $x(\nu_n)\ge x_n(D,\eta)$.
\end{proposition}

\begin{proof}
At $n=0$ the two laws are equal. Suppose $\nu_n$ is more informative than $\mu_n=\mu_n(D,\eta)$. By
\cref{lem:recursion}, $\mu_{n+1}=(E_\eta\mu_n)^{\circledast D}$, and since fusion is associative and
commutative on posterior laws, the right side may be grouped as in either scheme. In that grouping replace each
factor $E_\eta\mu_n$ by $E_\eta\nu_n$, one at a time, and apply $\Pi$ where the scheme applies it. Each
replacement and each application of $\Pi$ can only make the law more informative, so $\nu_{n+1}$ is more
informative than $\mu_{n+1}$. The last assertion follows from the first, as in \cref{lem:compare}.
\end{proof}

No closeness of $\nu_n$ to $\mu_n$ is used: every step can only add information, and the inequality holds however
large the losses from the grid and the rounding are. When two laws on $G_K$ are fused, pairs with
$\langle p,r\rangle=0$ have probability zero and are dropped before normalizing.

\subsection{Exact arithmetic}
Masses are stored as integers summing to $S$, and every transition between orbits as a vector of nonnegative
integers with a common denominator. A grid point $v/K$ shrunk by $E_\eta$ has coordinates
$((B-A)K+3Av_i)/(3BK)$, so the split of a single shrunken point has denominator $3B$. For a fusion, fix one
orientation of the left orbit and average over the $o\in\{1,3,6\}$ orientations of the right orbit; by
simultaneous permutation of the colours this gives the full grouped transition. If the likelihood products
are $z_i$ with $C=z_1+z_2+z_3$, the split weights of $z/C$ have denominator $C$, the pair has mass proportional
to $C$, and $C$ cancels. With the factor $36/o$, which is an integer, the transitions become integer vectors with
row sum $108K^2$ for the fusion of two grid laws, $108(BK)^2$ for the first stage of sequential fusion and
$108BK^2$ for its later stages.

A new orbit mass is the integer quotient of an accumulator $\sum m_im_jt_{ij}$ by $S$ times the row sum. All
summands are nonnegative, so each accumulator is at most $S^2$ times the row sum. This is below $2^{128}$ in
every row of \cref{prop:entry} except $d=2$, where the computation uses arbitrary-precision integers; the
single-law step for binary fusion stays below $S\cdot3B<2^{64}$. The integer quotient is exactly the rounding of
\cref{lem:round}, and the removed units go to the vertex orbit. Finally
\[
x(\nu_n)=\frac1S\sum_{O}m_O\,\frac{3|v_O|^2-K^2}{3K^2},
\]
where $v_O\in\mathbb Z^3$ is any point of the orbit $O$. No floating-point number enters a coefficient, a mass
or a bound. The tables are checked as described in \cref{sec:computations}.

\subsection{The entry bounds}

\begin{proposition}\label{prop:entry}
For each row of \cref{tab:entry}, each degree $d$ in the range and each $n\ge n_0$,
\[
x_n\big(d,d^{-1/2}\big)\le x(\nu_{n_0})<r_d,
\]
where $\nu_{n_0}$ is computed with the listed $D$ and $\eta$, $K=96$ and $S=2^b$.
\end{proposition}

{\small
\begin{longtable}{lllllll}
\caption{Entry bounds. The column $x(\nu_{n_0})$ is the exact rational value rounded up to five places, and the
last column is $r_d$ at the smallest degree of the range.}\label{tab:entry}\\
\hline
$d$ & $D$ & $\eta$ & $n_0$ & $b$ & $x(\nu_{n_0})\le$ & $r_d$\\
\hline
\endfirsthead
\hline
$d$ & $D$ & $\eta$ & $n_0$ & $b$ & $x(\nu_{n_0})\le$ & $r_d$\\
\hline
\endhead
2 & 2 & $707107/1000000$ & 500 & 40 & 0.01142 & 0.01150\\
3 & 3 & $2887/5000$ & 106 & 40 & 0.01847 & 0.03000\\
4 & 4 & $1/2$ & 82 & 40 & 0.02108 & 0.03472\\
5 & 5 & $161/360$ & 25 & 32 & 0.04689 & 0.04822\\
6 & 6 & $89/218$ & 26 & 32 & 0.04765 & 0.04863\\
7 & 7 & $127/336$ & 22 & 32 & 0.05348 & 0.05374\\
8 & 8 & $99/280$ & 20 & 32 & 0.05732 & 0.05804\\
9 & 9 & $1/3$ & 18 & 32 & 0.06158 & 0.06173\\
10 & 10 & $37/117$ & 17 & 32 & 0.06472 & 0.06494\\
11 & 11 & $79/262$ & 17 & 32 & 0.06510 & 0.06776\\
12 & 12 & $97/336$ & 16 & 32 & 0.06822 & 0.07028\\
13 & 13 & $38/137$ & 15 & 32 & 0.07172 & 0.07253\\
14 & 14 & $89/333$ & 15 & 32 & 0.07205 & 0.07457\\
15 & 15 & $55/213$ & 14 & 32 & 0.07595 & 0.07643\\
16 & 16 & $1/4$ & 14 & 32 & 0.07601 & 0.07812\\
17 & 17 & $65/268$ & 14 & 32 & 0.07676 & 0.07969\\
18 & 18 & $33/140$ & 13 & 32 & 0.08104 & 0.08113\\
19 & 19 & $53/231$ & 13 & 32 & 0.08147 & 0.08247\\
20--21 & 21 & $53/237$ & 21 & 32 & 0.08481 & 0.08567\\
22--23 & 23 & $71/333$ & 24 & 32 & 0.08552 & 0.08598\\
24--25 & 25 & $69/338$ & 21 & 32 & 0.08746 & 0.08797\\
26--27 & 27 & $61/311$ & 19 & 32 & 0.08937 & 0.08975\\
28--29 & 29 & $55/291$ & 18 & 32 & 0.09042 & 0.09135\\
30--31 & 31 & $65/356$ & 17 & 32 & 0.09180 & 0.09280\\
32--33 & 33 & $32/181$ & 16 & 32 & 0.09384 & 0.09412\\
34--35 & 35 & $59/344$ & 16 & 32 & 0.09364 & 0.09534\\
36--37 & 37 & $1/6$ & 15 & 32 & 0.09570 & 0.09645\\
38--39 & 39 & $55/339$ & 15 & 32 & 0.09609 & 0.09748\\
40--41 & 41 & $37/234$ & 15 & 32 & 0.09607 & 0.09844\\
42--43 & 43 & $25/162$ & 14 & 32 & 0.09904 & 0.09933\\
44--45 & 45 & $49/325$ & 14 & 32 & 0.09922 & 0.10017\\
46--47 & 47 & $32/217$ & 14 & 32 & 0.09948 & 0.10095\\
48--49 & 49 & $41/284$ & 14 & 32 & 0.09976 & 0.10169\\
50--52 & 52 & $43/304$ & 19 & 32 & 0.10125 & 0.10238\\
53--55 & 55 & $25/182$ & 18 & 32 & 0.10238 & 0.10335\\
56--58 & 58 & $29/217$ & 17 & 32 & 0.10398 & 0.10425\\
59--61 & 61 & $25/192$ & 17 & 32 & 0.10418 & 0.10508\\
62--64 & 64 & $39/307$ & 17 & 32 & 0.10448 & 0.10585\\
65--67 & 67 & $33/266$ & 16 & 32 & 0.10634 & 0.10657\\
68--70 & 70 & $41/338$ & 16 & 32 & 0.10678 & 0.10725\\
71--73 & 73 & $40/337$ & 16 & 32 & 0.10694 & 0.10788\\
74--76 & 76 & $5/43$ & 16 & 32 & 0.10738 & 0.10847\\
77--79 & 79 & $31/272$ & 16 & 32 & 0.10758 & 0.10904\\
80--83 & 83 & $35/313$ & 66 & 32 & 0.10487 & 0.11933\\
84--88 & 88 & $37/339$ & 19 & 32 & 0.11856 & 0.11933\\
89--90 & 90 & $529999/5000000$ & 11 & 32 & 0.11761 & 0.11933\\
\hline
\end{longtable}}

\begin{proof}
In each row $D$ is the largest degree of the range and $\eta^2d\ge1$ at its smallest degree, so
$\eta\ge d^{-1/2}$ for every $d$ in the range. By \cref{lem:tail}, \cref{lem:compare} and
\cref{prop:dominate}, for $n\ge n_0$,
\[
x_n(d,d^{-1/2})\le x_{n_0}(d,d^{-1/2})\le x_{n_0}(D,\eta)\le x(\nu_{n_0}).
\]
The inequality $x(\nu_{n_0})<r_d$ is checked in exact arithmetic at every degree of the range; against the
curve in \eqref{eq:radii} it is checked by squaring. Degrees $2\le d\le5$ use sequential fusion and $d\ge6$
binary fusion.
\end{proof}

The rows are tight in places. At $d=2$ the bound is $0.011415$ against $r_2=0.0115$; at $d=9$ and $d=18$ the
margins are about $1.5\times10^{-4}$ and $1.0\times10^{-4}$.

\section{Large degrees}\label{sec:large}

For $d\ge91$ we do not compute an entry bound. Instead we show that a fixed point with $T\ge0$ cannot have
second moment $x\ge\tfrac1{10}$; together with \cref{prop:local3} this excludes every nontrivial fixed point. The
argument compares one step of the recursion with the Gaussian map of Sly~\cite[(4.7)]{Sly11}, by replacing the
branches one at a time, in the manner of the Lindeberg argument of \cite[Lemma~4.3]{Sly11} and
\cite[Appendix~A]{MSS22}. The new points are explicit constants, and a cancellation of the third-order terms that
uses a second observable.

\subsection{The Gaussian experiment}
Let $P=I-\tfrac13\one\one^{\mathsf T}$ and let $F(z)=\mathrm{softmax}(z)$, so $F_i(z)=e^{z_i}/\sum_je^{z_j}$. For
$s\ge0$ let $Z$ be a random vector which, given $\sigma=i$, is Gaussian with mean $\tfrac{9s}2Pe_i$ and covariance
$\tfrac{9s}2P$, with $\sigma$ uniform. The means have equal norms in the metric of the covariance, so the
likelihood ratios are $e^{Z_i-Z_j}$ and $F(Z)$ is the posterior of $\sigma$ given $Z$. Put
\[
g_G(s)=\E\big[F_1(Z)\mid\sigma=1\big]-\tfrac13,\qquad U_G(s)=\E\big[|3F(Z)-\one|^4\mid\sigma=1\big].
\]
By calibration $g_G(s)$ is the second moment $x$ of the posterior law of $Z$. This is the map of
\cite[(4.7)]{Sly11} for $q=3$ in our normalization: there the mean is $s(\tfrac32,-3,-3)$ before centring and the
covariance has diagonal $3$ and off-diagonal $-\tfrac32$, which after adding $\tfrac{3s}2$ to every logit is the law
above.

\begin{lemma}\label{lem:gap}
$s-g_G(s)\ge\max\{\tfrac9{20}s^2,\tfrac{13}{200}s\}>0$ for $\tfrac1{10}\le s\le\tfrac23$.
\end{lemma}

\begin{proof}
For $0<a\le b$ the experiment at $a$ is a garbling of the one at $b$: if $Y_b$ has the law at $b$, then
$\sqrt{a/b}\,Y_b+\sqrt{1-a/b}\,Y'$, with $Y'$ an independent centred copy of the noise, has the law at $a$. Hence
$g_G$ is nondecreasing, as in \cref{lem:compare}, and for $l\le s\le r$
\begin{equation}\label{eq:gapmono}
s-g_G(s)\ge l-g_G(r).
\end{equation}

In coordinates, with $\xi=(\xi_1,\xi_2)$ standard Gaussian,
\[
g_G(s)=\E f_s(\xi)-\tfrac13,\qquad
f_s(\xi)=\Big[1+2\exp\big(-\tfrac{9s}2-\tfrac32\sqrt{3s}\,\xi_1\big)\cosh\big(\tfrac32\sqrt s\,\xi_2\big)\Big]^{-1},
\]
since the two columns $(\sqrt{2/3},-1/\sqrt6,-1/\sqrt6)$ and $(0,1/\sqrt2,-1/\sqrt2)$ are an orthonormal basis of the
plane $\one^\perp$. For $p=F(z)$ and a direction $h$ with $R=\max_ih_i-\min_ih_i$,
\[
D^2F_1(z)[h,h]=p_1\big((h_1-\textstyle\sum_jp_jh_j)^2-\mathrm{Var}_p(h)\big).
\]
The first term lies between $0$ and $p_1(1-p_1)^2R^2\le\tfrac4{27}R^2$, because
$\tfrac4{27}-p(1-p)^2=(p-\tfrac13)^2(\tfrac43-p)$, and the variance lies between $0$ and $R^2/4$. Hence
$|D^2F_1[h,h]|\le R^2/4\le|h|^2/2$, and the Hessian of $f_s$ has operator norm at most $\tfrac{9s}4$ everywhere.

Divide $[-6,6]$ into $96$ intervals of length $\tfrac18$. For an interval $I_j$ let $w_j$ be its Gaussian mass,
$m_j$ the conditional mean and $v_j$ the mass times the conditional variance, and put $P=\sum_jw_j$,
$V=\sum_jv_j$. On a product cell $I_j\times I_k$ expand $f_s$ about $(m_j,m_k)$: the linear term integrates to
zero and the remainder is at most $\tfrac{9s}8|\xi-(m_j,m_k)|^2$. Summing over the cells, whose integrated squared
distances add up to $2PV$, and using $0\le f_s\le1$ off the square,
\begin{equation}\label{eq:gapquad}
g_G(s)\le\sum_{j,k}w_jw_kf_s(m_j,m_k)+\tfrac{9s}4PV+1-P^2-\tfrac13.
\end{equation}
Let $t_i=10/(100-i)$ for $0\le i\le85$, so that $t_0=\tfrac1{10}$ and $t_{85}=\tfrac23$. In
\cref{sec:computations} the right side of \eqref{eq:gapquad} at $s=t_{i+1}$ is enclosed by a rational number
$U_i$, and exact arithmetic gives
\[
t_i-U_i>\tfrac9{20}t_{i+1}^2,\qquad t_i-U_i>\tfrac{13}{200}t_{i+1}\qquad(0\le i<85).
\]
By \eqref{eq:gapmono} with $l=t_i$ and $r=t_{i+1}$, both bounds hold on $[t_i,t_{i+1}]$.
\end{proof}

\subsection{One replacement}
Let $\mu$ be a symmetric posterior law with moments $M,U,T$ as in \cref{sec:local}, and $a=3p-\one$, $A=|a|^2$. Given
the colour of the parent, the logarithmic likelihood vector of one branch through an edge of strength $\lambda$ is
$P\log(1+\lambda a)$, computed coordinatewise; its Gaussian counterpart is $G\sim N(\lambda^2Mv_i,
\tfrac{\lambda^2M}2P)$ given $\sigma=i$, where $v_i=\tfrac12(e_i-\tfrac13\one)$. For a smooth function $f$ on
$\R^3$, invariant under adding a constant to every coordinate, suppose that at every point
\begin{equation}\label{eq:derivs}
\begin{gathered}
|Df[v_i]|\le g_1,\quad |D^2f[h,h]|\le c_2|Ph|^2,\quad |D^3f[h,h,h]|\le c_3|Ph|^3,\\
|D^3f[r,h,h]|\le c_{31}|Pr||Ph|^2,\quad |D^4f[h,h,h,h]|\le c_4|Ph|^4 .
\end{gathered}
\end{equation}
For $0<L<1$ put
\[
B_4=\tfrac{25}{288}+\tfrac{49}{100}\tfrac L{1-L},\qquad B_3=\tfrac16+\tfrac{49}{20}LB_4,\qquad
B_1=\tfrac5{24}+\tfrac{49}{20}LB_3,
\]
\begin{align*}
\Gamma_U(f;L)&=g_1\Big(\tfrac16+\tfrac{49}{20}LB_4\Big)+\tfrac{c_2}2\Big(\tfrac1{24}+B_3\big(\tfrac1{1-L}+1+\tfrac L2\big)+\tfrac16\Big)
+\tfrac{c_{31}B_1}{2(1-L)^2}+\tfrac{c_3}6+\tfrac{c_4}{24(1-L)^4},\\
\Gamma_M(f;L)&=\tfrac{c_2}{12}+\tfrac{5c_{31}}{24}+\tfrac{5c_3L^2}{72}+\tfrac{c_4(2+4L^2+L^4)}{24},
\end{align*}
and, for $\{i,j,k\}=\{1,2,3\}$,
\begin{equation}\label{eq:ell}
\ell_i^f=\tfrac14\Big(\textstyle\sum_rD_r^3f-D_if-(D_j-D_k)^2f\Big).
\end{equation}
Here $D_r=\partial/\partial z_r$, so that $D_r^3f=D^3f[e_r,e_r,e_r]$ and $(D_j-D_k)^2f=D^2f[e_j-e_k,e_j-e_k]$.

\begin{lemma}\label{lem:replace}
Let $0<\lambda\le L<1$ and let $W$ be any random vector, independent of the branch given $\sigma$. Then, given
$\sigma=i$,
\[
\E f\big(W+P\log(1+\lambda a)\big)-\E f(W+G)
\le\lambda^3T\,\E\,\ell_i^f(W)+\lambda^4\big(\Gamma_U(f;L)\,U+\Gamma_M(f;L)\,M^2\big).
\]
\end{lemma}

\begin{proof}
Take $i=1$; the other cases follow by permuting colours. Given $\sigma=1$, the child's point $a$ has law
$(1+\lambda a_1)\mu(da)$. Squares and cubes of vectors below are coordinatewise. Write
\[
\Lambda=P\log(1+\lambda a)=\lambda a+\lambda^2b+\lambda^3c+r_4,\qquad b=-\tfrac12Pa^2,\quad c=\tfrac13Pa^3.
\]
Since $a_1+a_2+a_3=0$ we have $|b|=A/(2\sqrt6)$, $c=Aa/6$, $|Pa^4|^2\le\tfrac{25}{216}A^4$ and $A\le6$. The
integral form of the remainder in $\log(1+t)=t-\tfrac{t^2}2+\tfrac{t^3}3-\dots$, valid for $t\ge-\lambda>-1$, then gives
\begin{equation}\label{eq:logrem}
\begin{gathered}
|r_4|\le\lambda^4A^2B_4,\qquad|\Lambda-\lambda a-\lambda^2b|\le\lambda^3A^{3/2}B_3,\qquad|\Lambda-\lambda a|\le\lambda^2AB_1,\\
|\Lambda|\le\frac{\lambda\sqrt A}{1-L},\qquad|\lambda a+\lambda^2b|\le\lambda\sqrt A\big(1+\tfrac L2\big).
\end{gathered}
\end{equation}
Let $m=\lambda^2Mv_1$ and $C=\tfrac{\lambda^2M}2P$, so that $G\sim N(m,C)$. Condition on $W$ and expand
$f(W+y)$ to fourth order in $y$, for $y=\Lambda$ and $y=G$; the fourth-order remainder is at most
$c_4|Py|^4/24$. We compare the two expansions term by term. For any equivariant plane-valued function $V$,
$\E_\mu[a_1V(a)]=v_1\E_\mu[a\cdot V(a)]$, and every radial quantity has the same mean under $\mu$ and under
$(1+\lambda a_1)\mu$.

\emph{First order.} By the identity just stated,
$\E_1\Lambda=m-\tfrac12\lambda^3Tv_1+\tfrac16\lambda^4Uv_1+\lambda v_1\E[a\cdot r_4]$. The term in $\lambda^3$ is
kept; by \eqref{eq:logrem} and $\sqrt6\le\tfrac{49}{20}$ the others contribute at most
$\lambda^4Ug_1(\tfrac16+\tfrac{49}{20}LB_4)$.

\emph{Second order.} The difference $\E_1\Lambda\Lambda^{\mathsf T}-\E GG^{\mathsf T}$ is
$-\tfrac12\lambda^3T(e_2-e_3)(e_2-e_3)^{\mathsf T}$ plus a matrix of nuclear norm at most
$\lambda^4\{U[\tfrac1{24}+B_3(\tfrac1{1-L}+1+\tfrac L2)+\tfrac16]+\tfrac16M^2\}$. Indeed,
$\|zz^{\mathsf T}-ww^{\mathsf T}\|_*\le|z-w|(|z|+|w|)$ compares $\Lambda$ with $\lambda a+\lambda^2b$; the term
$\lambda^4bb^{\mathsf T}$ costs $U/24$; the change of density acting on $\lambda^3(ab^{\mathsf T}+ba^{\mathsf T})$ costs
$U/6$; and $mm^{\mathsf T}$ has norm $\lambda^4M^2/6$. With the Taylor factor $\tfrac12$ and the bound $c_2$ on
the Hessian this gives the $c_2$ terms of $\Gamma_U$ and $\Gamma_M$.

\emph{Third order.} Let $q(y)=D^3f(W)[y,y,y]$. Along the segment from $\lambda a$ to $\Lambda$,
$|q(\Lambda)-q(\lambda a)|\le3c_{31}\lambda^4A^2B_1/(1-L)^2$. The change of density acting on $q(\lambda a)$ is at
most $c_3\lambda^4U$, because $|a_1|\le\sqrt A$, and what remains is $\lambda^3\E_\mu q(a)$, which carries the
cubic term. For the Gaussian, with $G=m+H$, $\E q(G)=q(m)+3\E D^3f(W)[m,H,H]$; here $\E|H|^2=\lambda^2M$,
$|m|=\lambda^2M/\sqrt6$, $M\le6$ and $1/\sqrt6\le\tfrac5{12}$, which bound the two terms by
$\lambda^4M^2\cdot\tfrac5{12}c_3L^2$ and $\lambda^4M^2\cdot\tfrac54c_{31}$. With the Taylor factor $\tfrac16$ this
gives the $c_3$ and $c_{31}$ terms.

\emph{Fourth order.} By \eqref{eq:logrem}, $\E_1|\Lambda|^4\le\lambda^4U/(1-L)^4$, and
$\E|G|^4=\lambda^4M^2(2+\tfrac23\lambda^2M+\tfrac1{36}\lambda^4M^2)\le\lambda^4M^2(2+4L^2+L^4)$.

Finally, the terms in $\lambda^3$ are $-\tfrac14TD_1f$ from the mean, $-\tfrac14T(D_2-D_3)^2f$ from the second
moment, and $\tfrac16\cdot\tfrac32T\sum_rD_r^3f$ from the third, since
$\E_\mu D^3f[a,a,a]=\tfrac32T\sum_rD_r^3f$ for a symmetric law. Their sum is $\lambda^3T\ell_1^f(W)$.
\end{proof}

\subsection{The two comparisons}
Write $Q(z)=|3F(z)-\one|^4$. Applied to $f=F_1$ and to $f=Q$, \cref{lem:replace} needs the constants in
\eqref{eq:derivs}.

\begin{lemma}\label{lem:softmax}
For $f=F_1$ the bounds \eqref{eq:derivs} hold with $(g_1,c_2,c_3,c_{31},c_4)=(\tfrac18,\tfrac15,\tfrac25,\tfrac12,\tfrac35)$,
and for $f=Q$ with $(8,20,38,38,99)$.
\end{lemma}

\begin{proof}
For $F_1$, $DF_1[v_i]=\tfrac12p_1(\delta_{1i}-p_i)$, which is at most $\tfrac18$ in absolute value. Both functions
are invariant under adding a constant to $z$, so every derivative sees only $Ph$, and a zero direction is trivial.
For the remaining bounds translate and rescale $h$ so that its smallest coordinate is $0$ and its largest is $1$.
Using $h\mapsto-h$ and the symmetry between colours $2$ and $3$, it suffices to take $h=(0,t,1)$ or $h=(t,0,1)$
with $0\le t\le1$ for $F_1$, and $h=(0,t,1)$ for the symmetric function $Q$; in all cases
$|Ph|^2=\tfrac23(1-t+t^2)>0$. Write $p=(u,(1-u)v,(1-u)(1-v))$ with $(u,v)\in[0,1]^2$, which covers the closed
simplex. With $m=\sum_ip_ih_i$, $r=h_1-m$ and the cumulants $k_2,k_3,k_4$ of $h$ under $p$, differentiation of
$p_i'=p_i(h_i-m)$ gives
\[
D^2F_1=p_1(r^2-k_2),\quad D^3F_1=p_1(r^3-3rk_2-k_3),\quad D^4F_1=p_1(r^4-6r^2k_2-4rk_3+3k_2^2-k_4),
\]
and the derivatives of $Q=(9\sum_ip_i^2-3)^2$ are obtained in the same way. Each bound is then a polynomial
inequality in $(u,v,t)\in[0,1]^3$ with rational coefficients: $\tfrac15|Ph|^2\pm D^2F_1\ge0$,
$\tfrac4{25}|Ph|^6-(D^3F_1)^2\ge0$, $\tfrac35|Ph|^4\pm D^4F_1\ge0$, and for $Q$ the inequalities
$8\pm\tfrac12D_1Q\ge0$, $20|Ph|^2\pm D^2Q\ge0$, $38^2|Ph|^6-(D^3Q)^2\ge0$, $99|Ph|^4\pm D^4Q\ge0$ and
$38^2|Ph|^4-\sum_j(V_j)^2\ge0$, where $V_j=D^3Q[e_j,h,h]$. These eighteen inequalities are proved in
\cref{sec:computations} by Bernstein expansion with exact subdivision. The last one gives the mixed bound for
$Q$, since $V\perp\one$ and $D^3Q[r,h,h]=\sum_jr_jV_j$. For $F_1$ the mixed bound follows from the cubic one by
polarization: for unit vectors $r,h$ in the plane,
\[
3D^3F_1[h,h,r]=D^3F_1[(h+\tfrac r2)^{3}]-D^3F_1[(h-\tfrac r2)^{3}]-\tfrac14D^3F_1[r^{3}],
\]
and $|h+\tfrac r2|^3+|h-\tfrac r2|^3\le\tfrac72$, so $|D^3F_1[h,h,r]|\le\tfrac13\cdot\tfrac25\cdot\tfrac{15}4=\tfrac12$.
\end{proof}

The third-order terms have a sign. At a point $z$ with $p=F(z)$, let $e_2=\sum_{i<j}p_ip_j$ and $e_3=p_1p_2p_3$, and
write $\ell_i^{F_i}$ for \eqref{eq:ell} applied to $f=F_i$.

\begin{lemma}\label{lem:cubic}
At every $z\in\R^3$, with $\beta=\tfrac{49}4$,
\[
\sum_ip_i\ell_i^{F_i}(z)=-e_2(e_2-9e_3)\le0,\qquad
\sum_ip_i\ell_i^{Q}(z)+9\beta\sum_ip_i\ell_i^{F_i}(z)\le0.
\]
\end{lemma}

\begin{proof}
Since $D_jF_i=p_i(\delta_{ij}-p_j)$, both sums are polynomials in $p$, and direct expansion gives
\[
\sum_ip_i\ell_i^{F_i}=-e_2(e_2-9e_3),\qquad
\sum_ip_i\ell_i^{Q}=54\big(18e_2^3-270e_2^2e_3-5e_2^2+99e_2e_3-6e_3\big).
\]
By Maclaurin's inequality $e_2\ge3e_3^{2/3}$, and $e_3\le\tfrac1{27}$, so $e_2\ge9e_3$; this is the first claim.

For the second, put $a=3p-\one$, $A=|a|^2$ and $J=\sum_ia_i^3$, so that $e_2=(6-A)/18$ and
$e_3=(J+3-\tfrac32A)/81$. Substitution shows that the left side of the second inequality equals $C(A)+B(A)J$, where
\[
C(A)=\big(4-\tfrac\beta3\big)A+\big(\tfrac\beta{18}-4\big)A^2+\tfrac23A^3,\qquad
B(A)=\tfrac\beta3-2+\big(3-\tfrac\beta{18}\big)A-\tfrac59A^2.
\]
On the simplex $0\le A\le6$. Since $a_1+a_2+a_3=0$, the discriminant identity
$A^3/6-J^2=\tfrac13\prod_{i<j}(a_i-a_j)^2$ gives $J\le A^{3/2}/\sqrt6$, and $e_3\ge0$ gives $J\ge\tfrac32A-3$.
If $B(A)\ge0$, put $t=\sqrt{A/6}\in[0,1]$ and use $J\le6t^3$:
\[
C+BJ\le-2t^2(1-t)^2\Big[48\big(t-\tfrac{23}{384}\big)^2+\tfrac{239}{3072}+60t^3\Big]\le0.
\]
If $B(A)<0$, use $J\ge\tfrac32A-3$:
\[
C+BJ\le-\tfrac1{144}(6-A)^2(24A+25)\le0.
\]
Both right sides are exact polynomial identities in $t$ and $A$ respectively.
\end{proof}

\begin{proposition}\label{prop:compare}
Let $\mu$ be a symmetric posterior law with $T(\mu)\ge0$, let $d\ge1$, $0<\lambda\le L<1$, and let $x',U'$ be the
second and fourth moments $x$ and $U$ of $T_\lambda\mu$. With $s=d\lambda^2x$, $\beta=\tfrac{49}4$,
$f_U=\Gamma_U(F_1;L)$, $f_M=\Gamma_M(F_1;L)$, $q_U=\Gamma_U(Q;L)$ and $q_M=\Gamma_M(Q;L)$,
\begin{align}
x'-g_G(s)&\le d\lambda^4(f_UU+f_MM^2),\label{eq:cmp1}\\
U'-U_G(s)+9\beta\big(x'-g_G(s)\big)&\le d\lambda^4\big((q_U+9\beta f_U)U+(q_M+9\beta f_M)M^2\big).\label{eq:cmp2}
\end{align}
\end{proposition}

\begin{proof}
Given $\sigma$, the logarithmic likelihood of the $d$ branches is the sum of $d$ independent copies of
$P\log(1+\lambda a)$, and that of the Gaussian experiment at $s=d\lambda^2x$ is the sum of $d$ independent copies of
$G$, since $d\lambda^2M/2=\tfrac{9s}2$. Replace the branches by Gaussian copies one at a time. Every intermediate
sum is the logarithmic likelihood of a genuine symmetric experiment, since sums of independent likelihoods remain
likelihoods, so its posterior is $F(W)$ and its law $\nu$ satisfies $\PP(W\in\cdot\mid\sigma=i)=3F_i(W)\,\nu$.
Hence the third-order term of \cref{lem:replace} for $f=F_\sigma$, averaged over $\sigma$, is $\lambda^3T$ times
the $\nu$-average of $\sum_ip_i\ell_i^{F_i}$, and for $f=Q$ it is $\lambda^3T$ times the $\nu$-average of
$\sum_ip_i\ell_i^Q$. By calibration $\E F_\sigma(\cdot)=x+\tfrac13$ at both ends of the chain, with $x=x'$ for the
branches and $x=g_G(s)$ for the Gaussian, and $Q$ is symmetric, so its conditional and unconditional means agree.
Summing the $d$ replacements, applying \cref{lem:replace} to $F_\sigma$ for \eqref{eq:cmp1} and to $Q$ and
$9\beta F_\sigma$ separately for \eqref{eq:cmp2}, the third-order terms are nonpositive by \cref{lem:cubic} and
$T\ge0$, and may be dropped.
\end{proof}

\subsection{Exclusion of large fixed points}

\begin{proposition}\label{prop:large}
Let $d\ge91$ and $\lambda=d^{-1/2}$. Every symmetric posterior law $\mu$ with $T_\lambda\mu=\mu$ and $T(\mu)\ge0$
has $x(\mu)<\tfrac1{10}$.
\end{proposition}

\begin{proof}
Put $D=91$, $L=209657/2000000$ and $\beta=\tfrac{49}4$. Then $DL^2\ge1$, so $\lambda\le L$, and
\cref{prop:compare} applies. Let $f_U,f_M,q_U,q_M$ be as there, at this $L$, and define
\[
E_d(s)=\frac{d-q_U}{f_U}\big(s-g_G(s)\big)-81\Big[\Big(1-\frac{q_U}d\Big)\frac{f_M}{f_U}+\frac{q_M}d\Big]s^2-U_G(s).
\]
In \cref{sec:computations} we verify, in exact rational arithmetic,
\begin{equation}\label{eq:d91}
D-q_U-9\beta f_U>0,\qquad q_M-q_U\frac{f_M}{f_U}>0,\qquad E_D(s)>\tfrac3{200}\quad\big(\tfrac1{10}\le s\le\tfrac23\big).
\end{equation}

Suppose $x=x(\mu)\ge\tfrac1{10}$; then $x\le\tfrac23$. At the fixed point $x'=x$, $U'=U$, $s=d\lambda^2x=x$ and
$d\lambda^4=1/d$. Write $M=9x$, $K=U/(81x^2)$ and $\gamma=(x-g_G(x))/x^2$. Then \eqref{eq:cmp1} reads
$K\ge(d\gamma/81-f_M)/f_U$, and \eqref{eq:cmp2} reads
\[
\big(d-q_U-9\beta f_U\big)K\le\frac{dU_G(x)}{81x^2}-\frac{\beta d\gamma}9+q_M+9\beta f_M.
\]
The coefficient on the left is positive by \eqref{eq:d91}, since $d\ge D$. Inserting the lower bound for $K$, the
terms in $\beta$ cancel, and after multiplication by $81x^2/d$ what remains is $E_d(x)\le0$.

By \cref{lem:gap}, $x-g_G(x)>0$. The coefficient $(d-q_U)/f_U$ increases with $d$, while
$81[(1-q_U/d)f_M/f_U+q_M/d]=81[f_M/f_U+(q_M-q_Uf_M/f_U)/d]$ decreases with $d$ by \eqref{eq:d91}. Hence
$E_d(x)\ge E_D(x)>\tfrac3{200}$, a contradiction.
\end{proof}

\begin{remark}\label{rem:ug}
Without the second observable, \eqref{eq:cmp1} alone does not bound $U$, and the third-order term of $Q$ by itself
has no sign: $\sum_ip_i\ell_i^Q>0$ at $p=(\tfrac25,\tfrac3{10},\tfrac3{10})$. The multiple $9\beta F_\sigma$ of the
first observable corrects this sign everywhere on the simplex (\cref{lem:cubic}), and its contribution cancels in
the proof of \cref{prop:large}. The value $\beta=\tfrac{49}4$ is one convenient choice; the same proof works for
any $\beta$ for which $C(A)+B(A)J\le0$ on the feasible region and \eqref{eq:d91} holds.
\end{remark}

\section{Proof of \texorpdfstring{\cref{thm:main}}{Theorem 1.1}}\label{sec:proof}

By \cref{lem:compare} with $D=d$, it suffices to take $\lambda=d^{-1/2}$. The ingredients for each degree are
collected in \cref{tab:degrees}.

\begin{table}[ht]
\caption{The two estimates used at each degree.}\label{tab:degrees}
\begin{tabular}{lll}
\hline
degrees & exclusion near $\delta_u$ & entry or exclusion away from $\delta_u$\\
\hline
$d=2$ & $x\le r_2=\tfrac{23}{2000}$, \cref{prop:local2} & $x_{500}\le0.011415$, \cref{prop:entry}\\
$d=3$ & $x\le r_3=\tfrac3{100}$, \cref{prop:local3} & $x_{106}\le0.018465$, \cref{prop:entry}\\
$4\le d\le90$ & $x\le r_d$ in \eqref{eq:radii}, \cref{prop:local3} & \cref{tab:entry}, \cref{prop:entry}\\
$d\ge91$ & $x\le r_d$, $r_d\ge\tfrac{179}{1500}$, \cref{prop:local3} & $x(\mu)<\tfrac1{10}$ if $T(\mu)\ge0$, \cref{prop:large}\\
\hline
\end{tabular}
\end{table}

\begin{proof}[Proof of \cref{thm:main}]
Let $\lambda=d^{-1/2}$.

\emph{Degrees $2\le d\le90$.} By \cref{prop:entry} there is $N$ with $x_N<r_d$. By \cref{prop:local2} for $d=2$
and \cref{prop:local3} for $d\ge3$, every symmetric posterior law $\mu$ with $T_\lambda\mu=\mu$ and
$x(\mu)\le r_d$ equals $\delta_u$. \Cref{prop:reduction} with $r=r_d$ gives nonreconstruction.

\emph{Degrees $d\ge91$.} By \cref{lem:tail}, $\mu_\infty$ is a symmetric posterior law with
$T_\lambda\mu_\infty=\mu_\infty$. Since $d\lambda^3=d^{-1/2}<1$, \cref{cor:sign} gives $T(\mu_\infty)\ge0$, and
\cref{prop:large} gives $x(\mu_\infty)<\tfrac1{10}$. For $d\ge80$ we have $r_d\ge\varrho_d=\tfrac{179}{1500}>\tfrac1{10}$,
so $\mu_\infty=\delta_u$ by \cref{prop:local3}. Hence $x_n\to0$, and nonreconstruction follows as in the proof of
\cref{prop:reduction}.

Finally, for $\lambda<d^{-1/2}$, \cref{lem:compare} gives $x_n(d,\lambda)\le x_n(d,d^{-1/2})\to0$. Since
reconstruction holds for $\lambda>d^{-1/2}$ by~\cite{KS66}, the threshold is $d^{-1/2}$.
\end{proof}

\section{The computations}\label{sec:computations}

Every inequality used above that is not proved by hand is a finite list of inequalities between rational
numbers. This section describes how each list is produced and checked. The programs are in the accompanying
repository; each check has a producer and a separate replay that does not import the producer.

\subsection{Bernstein certificates}
\Cref{lem:moments,prop:local2,prop:local3} and \cref{lem:softmax} reduce to the nonnegativity of polynomials with
rational coefficients on a box, which after an affine change of variables is $[0,1]^k$ with $2\le k\le4$. Where
a polynomial vanishes on part of the boundary, the monomial factors responsible are divided out first, exactly,
and the quotient is treated. For a polynomial $\sum_ea_ex^e$ of degree $n_j$ in $x_j$ the
Bernstein coefficients on the box are
\[
b_\kappa=\sum_{e\le\kappa}a_e\prod_j\binom{\kappa_j}{e_j}\Big/\binom{n_j}{e_j},
\]
and the polynomial is a convex combination of them at every point of the box. If some coefficient is negative,
the box is bisected, the coordinates taken in turn, and the coefficients on the two halves are obtained exactly by
de Casteljau's algorithm. A certificate is a prefix-free set of dyadic subboxes of total volume one, each with
nonnegative coefficients; a negative coefficient is never taken as a counterexample. All arithmetic is in
integers with a common denominator. For \cref{lem:softmax} the ten inequalities for $F_1$ close with $68$ leaves
and the eight for $Q$ with $124$ leaves. Exact negative controls show that the constants cannot be lowered to
$\tfrac1{10}$, $\tfrac15$ and $\tfrac25$ in place of $\tfrac15$, $\tfrac25$ and $\tfrac35$: at $p=(\tfrac14,0,\tfrac34)$ and
$h=(0,\tfrac12,1)$ one has $|Ph|^2=\tfrac12$, $|D^2F_1|=\tfrac3{32}$ and $|D^4F_1|=\tfrac{15}{128}$, and at
$p=(\tfrac12,0,\tfrac12)$ with the same $h$, $|D^3F_1|=\tfrac18$.

\subsection{Entry bounds}
The laws $\nu_n$ of \cref{sec:envelope} are computed with integer masses, as described there, from the parameters
in \cref{tab:entry}. The output of each row is the exact rational number $x(\nu_{n_0})$, stored as a fraction;
the comparison with $r_d$ is made in exact arithmetic for every $d$ in the row, by squaring where $r_d$ involves
$\sqrt d$. The computation is deterministic, and a replay recomputes each row from the parameters alone.

\subsection{Gaussian quadrature}
\Cref{lem:gap} and the third inequality in \eqref{eq:d91} need upper bounds for $g_G$ and $U_G$ at the points
$t_i=10/(100-i)$, $0\le i\le85$. For \cref{lem:gap} they are obtained from \eqref{eq:gapquad}. For \eqref{eq:d91}
the same quadrature is used with $192$ intervals of length $\tfrac1{16}$ on $[-6,6]$, for $F_1$ and for $Q$; for $Q$ the cell error is at most $90s\,PV$, since $|D^2Q[h,h]|\le20|Ph|^2$ and the map from $\xi$ to the
logits has squared norm $\tfrac{9s}2$, and the tail is at most $36(1-P^2)$, since $0\le Q\le36$. For $F_1$ the
same argument with $c_2=\tfrac15$ gives the cell error $\tfrac{9s}{10}PV$. The masses $w_j$, the moments $m_j$
and $v_j$, the exponentials and the sums are evaluated in ball arithmetic with $160$ bits~\cite{Joh17}, with
outward rounding throughout, and the upper endpoints of the resulting balls are exported as exact dyadic
rationals. All further comparisons use these rationals. A replay recomputes the endpoint bounds from the stored
cell data with independent code and, at $t_0$, $t_{42}$ and $t_{85}$, evaluates the quadrature sums afresh on all
$192^2$ cells and checks that the two enclosures overlap.

\subsection{The large-degree certificate}
At $L=\tfrac{209657}{2000000}$ the constants of \cref{prop:compare} are
\[
f_U=0.27741\ldots,\quad f_M=0.17224\ldots,\quad q_U=27.0695\ldots,\quad q_M=18.0441\ldots,
\]
all exact rationals. Then $D-q_U-9\beta f_U=33.346\ldots$ and $q_M-q_Uf_M/f_U=1.2367\ldots$, which are the first
two inequalities of \eqref{eq:d91}. For the third write
$E_D(s)=\alpha s-\alpha g_G(s)-Bs^2-U_G(s)$ with $\alpha=(D-q_U)/f_U=230.458\ldots$ and
$B=81[(1-q_U/D)f_M/f_U+q_M/D]=51.393\ldots$. The upper bounds for $g_G$ and $U_G$ give lower bounds for $E_D$ at
the $86$ points $t_i$, and the following lemma controls $E_D$ between them.

\begin{lemma}\label{lem:curv}
For $s>0$, $|g_G''(s)|\le\tfrac{1269}{40}<32$ and $|U_G''(s)|\le\tfrac{14229}4<3600$.
\end{lemma}

\begin{proof}
Given $\sigma=1$, the vector $Z(s)$ is a Brownian motion in the plane $\one^\perp$ with drift
$m_0=(3,-\tfrac32,-\tfrac32)$ and covariance $\tfrac92P$ per unit time, so for a smooth function $f$ with
bounded derivatives $\tfrac{d^2}{ds^2}\E f(Z(s))=\E\,\mathcal L^2f(Z(s))$ with
$\mathcal L=m_0\cdot\nabla+\tfrac94\Delta_P$. Let $e,e'$ be an orthonormal basis of the plane. For a symmetric
$4$-tensor with $|D^4f[h^4]|\le c_4|h|^4$, the identity
$12D^4f[e,e,e',e']=D^4f[(e+e')^4]+D^4f[(e-e')^4]-2D^4f[e^4]-2D^4f[e'^4]$ gives $|D^4f[e,e,e',e']|\le c_4$, and so
$|\Delta_P^2f|\le4c_4$. With $|m_0|^2=\tfrac{27}2$ and $|m_0|<\tfrac{15}4$,
\[
|\mathcal L^2f|\le\tfrac{27}2c_2+2\cdot\tfrac94\cdot2c_{31}|m_0|+\tfrac{81}{16}\cdot4c_4
\le\tfrac{27}2c_2+\tfrac{135}4c_{31}+\tfrac{81}4c_4 .
\]
The constants of \cref{lem:softmax} give $\tfrac{1269}{40}$ for $f=F_1$ and $\tfrac{14229}4$ for $f=Q$.
\end{proof}

Hence $|E_D''|\le C_D=32\alpha+2B+3600=11077.4\ldots$, and on each interval $[t_i,t_{i+1}]$
\[
E_D(s)\ge\min\{E_D(t_i),E_D(t_{i+1})\}-\tfrac18C_D(t_{i+1}-t_i)^2 .
\]
For all $85$ intervals the right side, computed from the exported lower bounds, exceeds $0.01584>\tfrac3{200}$.
Two controls are part of the check: raising the bound for $U_G$ by one at every point makes it fail, and so
does removing one endpoint.

\begin{thebibliography}{99}
\bibitem{BRZ95} P.~M. Bleher, J.~Ruiz and V.~A. Zagrebnov, On the purity of the limiting Gibbs state for
the Ising model on the Bethe lattice, \emph{J. Stat. Phys.} 79 (1995), 473--482.
\bibitem{BM11} N.~Bhatnagar and E.~Maneva, A computational method for bounding the probability of
reconstruction on trees, \emph{SIAM J. Discrete Math.} 25 (2011), 854--871; arXiv:0903.4812.
\bibitem{EKPS00} W.~Evans, C.~Kenyon, Y.~Peres and L.~J. Schulman, Broadcasting on trees and the Ising
model, \emph{Ann. Appl. Probab.} 10 (2000), 410--433.
\bibitem{Joh17} F.~Johansson, Arb: efficient arbitrary-precision midpoint-radius interval arithmetic,
\emph{IEEE Trans. Comput.} 66 (2017), 1281--1292.
\bibitem{KS66} H.~Kesten and B.~P. Stigum, Additional limit theorems for indecomposable multidimensional
Galton--Watson processes, \emph{Ann. Math. Statist.} 37 (1966), 1463--1481.
\bibitem{Mos01} E.~Mossel, Reconstruction on trees: beating the second eigenvalue, \emph{Ann. Appl.
Probab.} 11 (2001), 285--300.
\bibitem{MSS22} E.~Mossel, A.~Sly and Y.~Sohn, Exact phase transitions for stochastic block models and
reconstruction on trees, \emph{Ann. Probab.} (2025); arXiv:2212.03362.
\bibitem{Sly11} A.~Sly, Reconstruction for the Potts model, \emph{Ann. Probab.} 39 (2011), 1365--1406;
arXiv:0811.1208.
\end{thebibliography}

\end{document}
